\uuid{2vTY}
\titre{Standardiser une loi normale}
\auteur{Maxime Nguyen}
\organisation{AMSCC}
\datecreate{2026-11-24}
\niveau{L2}
\module{Probabilité}
\chapitre{Probabilité continue}
\sousChapitre{Loi normale}
\difficulte{1}
\theme{loi normale, standardisation, fonction de répartition}

\contenu{
\begin{question}
\qtype{onlyone}
On note $\mathcal{N}(\mu,\sigma)$ la loi normale d'espérance $\mu$ et d'écart-type $\sigma$, et $\Phi$ la fonction de répartition de $\mathcal{N}(0,1)$ : $\Phi(t) = \mathbb{P}(Z \leq t)$ pour $Z \sim \mathcal{N}(0,1)$. Soit $X \sim \mathcal{N}(10, 2)$. $\mathbb{P}(X \leq 13)$ est égale à :
\begin{answers}
\good{$\Phi(1{,}5)$}
\bad{$\Phi(13)$}
\bad{$\Phi(0{,}75)$}
\bad{$\Phi(6{,}5)$}
\end{answers}
\begin{explanations}
$\mathbb{P}(X \leq 13) = \mathbb{P}\left(\frac{X - 10}{2} \leq \frac{13 - 10}{2}\right) = \Phi(1{,}5) \approx 0{,}9332$.

Contre-sens visés : utiliser $\Phi$ sans standardiser (B) ; diviser par la
variance (C) ; oublier de centrer (D).
\end{explanations}
\end{question}
}
